\documentclass[10pt]{article}
\providecommand{\FIGDIR}{fig}
\newcommand{\DIGESTDATE}{3 September 2026}
% ===================== PM Research Watch - shared preamble =====================
\usepackage[T1]{fontenc}
\usepackage[utf8]{inputenc}
\usepackage{amsmath,amssymb}
\usepackage{tgpagella}
\usepackage{tgheros}
\renewcommand{\sfdefault}{qhv}
\renewcommand{\ttdefault}{lmtt}
\usepackage{microtype}
\usepackage[english]{babel}

\usepackage[a4paper,top=17mm,bottom=17mm,left=17mm,right=17mm,
            headsep=6mm,footskip=10mm]{geometry}
\usepackage{graphicx}
\usepackage{xcolor}
\usepackage{colortbl}
\usepackage{tikz}
\usetikzlibrary{calc}
\usepackage[most]{tcolorbox}
\usepackage{titlesec}
\usepackage{fancyhdr}
\usepackage{booktabs}
\usepackage{array}
\usepackage{longtable}
\usepackage{enumitem}
\usepackage{multicol}
\usepackage{ragged2e}
\usepackage[hidelinks]{hyperref}

% ------------------------------------------------------------------ palette
\definecolor{Ink}    {HTML}{14262E}
\definecolor{Deep}   {HTML}{0F4C5C}
\definecolor{Teal}   {HTML}{1F7A8C}
\definecolor{Sky}    {HTML}{4FA3C4}
\definecolor{Sage}   {HTML}{6FA287}
\definecolor{Amber}  {HTML}{E3A008}
\definecolor{Clay}   {HTML}{B5643C}
\definecolor{Coral}  {HTML}{D1495B}
\definecolor{Violet} {HTML}{6A5B9E}
\definecolor{Slate}  {HTML}{4A6572}
\definecolor{Paper}  {HTML}{FBF9F5}
\definecolor{Mist}   {HTML}{EEEAE2}
\definecolor{Rule}   {HTML}{D9D3C9}

\pagecolor{Paper}
\color{Ink}

\hypersetup{colorlinks=true, linkcolor=Deep, urlcolor=Teal, citecolor=Teal,
            pdftitle={PM Research Watch}, pdfauthor={Automated daily digest}}

% ------------------------------------------------------------------ headings
\titleformat{\section}
  {\sffamily\bfseries\large\color{Deep}}{}{0pt}
  {}[{\vspace{1.2mm}\color{Amber}\titlerule[1.1pt]}]
\titlespacing*{\section}{0pt}{5.5mm}{2.8mm}

\titleformat{\subsection}
  {\sffamily\bfseries\normalsize\color{Teal}}{}{0pt}{}
\titlespacing*{\subsection}{0pt}{3.6mm}{1.4mm}

\setlist[itemize]{leftmargin=4.2mm, itemsep=1.1mm, topsep=1.3mm, parsep=0pt,
                  label=\textcolor{Teal}{\footnotesize$\blacktriangleright$}}

% ------------------------------------------------------------------ header/footer
\pagestyle{fancy}
\fancyhf{}
\renewcommand{\headrulewidth}{0pt}
\renewcommand{\footrulewidth}{0pt}
\fancyhead[L]{\sffamily\scriptsize\color{Slate}PM RESEARCH WATCH\;\textcolor{Rule}{\textbar}\;\DIGESTDATE}
\fancyhead[R]{\sffamily\scriptsize\color{Slate}Particulate matter monitoring \& health effects}
\fancyfoot[C]{\sffamily\scriptsize\color{Slate}\thepage}

% ------------------------------------------------------------------ components
% metric tile
\newcommand{\metric}[3]{%
  \begin{tcolorbox}[enhanced, colback=white, colframe=Rule, boxrule=0.5pt,
    arc=1.6mm, left=2.2mm, right=2.2mm, top=1.8mm, bottom=1.8mm,
    borderline west={1.7mm}{0pt}{#3}, nobeforeafter, valign=center,
    height=17mm, width=\linewidth]
    \RaggedRight
    {\sffamily\bfseries\fontsize{17}{17}\selectfont\color{#3}#1}\\[0.5mm]
    {\sffamily\fontsize{6.5}{7.4}\selectfont\color{Slate}#2}
  \end{tcolorbox}}

% pull-quote / key insight box
\newtcolorbox{keybox}[1][]{enhanced, colback=Mist, colframe=Mist,
  arc=1.6mm, left=3.2mm, right=3.2mm, top=2.4mm, bottom=2.4mm,
  borderline west={1.7mm}{0pt}{Amber}, #1}

% subtopic band
\newcommand{\band}[2]{%
  \vspace{3.2mm}\noindent
  \begin{tcolorbox}[enhanced, colback=#2, colframe=#2, arc=1.2mm,
    left=2.6mm, right=2.6mm, top=1.1mm, bottom=1.1mm, boxrule=0pt,
    width=\linewidth, nobeforeafter]
    {\sffamily\bfseries\small\color{white}#1}
  \end{tcolorbox}\vspace{1.6mm}}

% one paper entry: {short cite}{title}{journal}{doi}{body}
\newcommand{\paperentry}[5]{%
  \par\vspace{1.9mm}\noindent
  {\sffamily\bfseries\footnotesize\color{Ink}#1}\;%
  {\sffamily\fontsize{7}{8}\selectfont\color{Slate}#3}\\[0.4mm]
  {\footnotesize\itshape\color{Deep}#2}\\[0.6mm]
  {\footnotesize #4}\\[0.3mm]
  {\fontsize{6.6}{7.6}\selectfont\color{Slate}\texttt{doi:}\,\href{https://doi.org/#5}{\texttt{\detokenize{#5}}}}
  \par}

% caption style
\newcommand{\figcap}[1]{\par\vspace{1.0mm}
  {\fontsize{7.2}{8.6}\selectfont\color{Slate}\textsf{\textbf{Fig.}}\; #1\par}\vspace{1.5mm}}

\providecommand{\permil}{\textperthousand}

\setlength{\parindent}{0pt}
\setlength{\parskip}{0pt}
\renewcommand{\arraystretch}{1.18}


\begin{document}

% ==================================================================== MASTHEAD
\thispagestyle{empty}
\begin{tikzpicture}[remember picture, overlay]
  \fill[Deep] (current page.north west) rectangle
      ($(current page.north east)+(0,-67mm)$);
  \fill[Amber] ($(current page.north west)+(0,-67mm)$) rectangle
      ($(current page.north east)+(0,-68.4mm)$);
  \foreach \x/\y/\r in {22/12/0.9, 41/26/1.5, 63/9/0.7, 88/31/1.1, 112/17/1.9,
                        138/29/0.8, 159/13/1.3, 178/54/1.0, 196/21/1.6,
                        30/44/0.6, 74/57/1.2, 125/48/0.9, 168/6/0.8, 191/59/0.7,
                        52/18/1.0, 100/60/0.6, 148/20/1.1, 66/49/1.4, 205/45/0.8}
    \fill[Sky, opacity=0.30] ($(current page.north west)+(\x mm,-\y mm)$) circle (\r mm);
\end{tikzpicture}

\vspace*{4mm}
{\sffamily\fontsize{7.6}{9}\selectfont\color{Sky}%
 AUTOMATED DAILY LITERATURE DIGEST\;\textbullet\;PhD RESEARCH WATCH\par}
\vspace{2.2mm}
{\sffamily\bfseries\fontsize{27}{29}\selectfont\color{white}%
 Particulate Matter\par}
{\sffamily\fontsize{27}{29}\selectfont\color{Sky}%
 Monitoring \& Health Effects\par}
\vspace{5.0mm}
{\sffamily\fontsize{8.4}{10}\selectfont\color{white}%
 Issue 2026-09-03\;\;\textcolor{Amber}{\textbar}\;\;Entry window 3 September 2026 (single day)\;\;%
 \textcolor{Amber}{\textbar}\;\;38 records in scope, 88 screened out\par}
\vspace{18mm}

% -------------------------------------------------------------- metric strip
\noindent
\begin{minipage}[t]{0.190\linewidth}\metric{38}{RECORDS IN SCOPE\\(88 SCREENED OUT)}{Deep}\end{minipage}\hfill
\begin{minipage}[t]{0.190\linewidth}\metric{9}{SUBTOPIC CLUSTERS\\OF TEN OCCUPIED}{Teal}\end{minipage}\hfill
\begin{minipage}[t]{0.190\linewidth}\metric{12}{EFFECT ESTIMATES\\(FOUR SOURCES)}{Amber}\end{minipage}\hfill
\begin{minipage}[t]{0.190\linewidth}\metric{11}{SENSING /\\INSTRUMENTATION}{Clay}\end{minipage}\hfill
\begin{minipage}[t]{0.190\linewidth}\metric{23}{RECORDS WITH NO\\HEALTH ENDPOINT}{Sky}\end{minipage}

\vspace{6mm}

% -------------------------------------------------------------- signal of day
\section*{\sffamily\bfseries\color{Deep}Signal of the day}
\vspace{-1.5mm}{\color{Amber}\rule{\linewidth}{1.1pt}}\vspace{2.5mm}

\begin{keybox}
{\sffamily\bfseries\small\color{Deep}Plantower's new PMSX003N is better than the PMS5003T
at dust and worse at humidity --- and the two failure modes are the same failure mode.}\\[1.4mm]
\footnotesize
Kaur, Mangin and Kelly (\textit{Aerosol Sci.\ Technol.}) co-located three PMSX003N, three
PMS5003T and three OPC-N3 against FEM monitors at a Utah DEQ site for five months
(Dec 2025--Apr 2026), through dust events, winter inversions, freezing fog, mist, snow and
rain. During dust the PMSX003N tracked the reference where its predecessor did not
(\textbf{PM$_{2.5}$ slope 1.85--1.91, $R^2>0.7$} against \textbf{slope 0.15, $R^2<0.37$}
for the PMS5003T; PM$_{10}$ \textbf{slope 0.91--1.07} against \textbf{$\sim$0.02}), and
under high humidity it \textbf{substantially overestimated both fractions}. Laboratory
size-binning explains both: the sensor \textbf{preferentially assigns particles $>$1\,$\mu$m
to the 2.5--5\,$\mu$m bin}. Coarse mineral dust is thereby recovered rather than lost, and
hygroscopically grown droplets are recruited into the same oversized bin. The design change
that fixes coarse-PM response is the design change that makes RH correction mandatory ---
so a humidity-uncorrected PMSX003N network is not a drop-in upgrade over a PMS5003T
network, it is a different bias. Note the PM$_{2.5}$ slope: even at its best the sensor
reads roughly \textbf{1.9$\times$} the reference during dust.
\end{keybox}

\vspace{3mm}

% -------------------------------------------------------------- exec brief
\section*{\sffamily\bfseries\color{Deep}What changed today}
\vspace{-1.5mm}{\color{Amber}\rule{\linewidth}{1.1pt}}\vspace{2.5mm}

\footnotesize
\begin{itemize}
\item \textbf{The largest single-day issue this watch has shipped, and it is instrument-led.}
38 records against a previous single-day maximum of 36 (27 August); only the retrospective
July back-scan is bigger. \textbf{Eleven} sit in \textit{Sensing, forecasting \&
instrumentation} and \textbf{twelve} in \textit{Exposure assessment \& modelling}: 23 of
38 before a single health endpoint. Volume is a one-day artefact of a fat \textit{ES\&T}
by-ISSN deposit, but the sensing half is substantive rather than incidental.

\item \textbf{Two independent low-cost-sensor papers, and they disagree about what LCS are for.}
Kaur et al.\ treat the sensor as an instrument to be characterised against a reference;
Guth et al.\ (\textit{Atmosphere}) treat it as the only practical way to get
\emph{time-resolved} emission factors across \textbf{19 prescribed burns}, where CO EFs
swing \textbf{+3260\%} with combustion phase. The second use case tolerates a bias the
first would reject, because within-burn contrast is the measurand.

\item \textbf{Component-resolved epidemiology produced two effect sizes that cannot both be
read the same way.} Chen K et al.\ report per-IQR frailty HRs of
\textbf{1.142 (1.055--1.236)} for organic matter and \textbf{1.131 (1.041--1.229)} for
black carbon in CHARLS ($n$=9{,}603) --- component ordering above bulk PM$_{2.5}$
(1.100), and biologically legible. Chen Y et al.\ report \textbf{OR 2.23 (2.19--2.27)}
for sulfate on congenital heart disease across \textbf{2.6 million} Guangdong births,
with the effect peaking in \emph{pregnancy week 1}. A doubling of CHD odds with a CI of
width 0.08 is not a plausible causal contrast; it is what a spatially smooth
satellite exposure does against a very large register. Cite the first, quarantine the second.

\item \textbf{A perspective worth reading against your own measurement plan.}
Gao et al.\ (\textit{ES\&T}) argue that the recurring \emph{inverse} O$_3$--health
associations are an artefact of measuring O$_3$ as a proxy for an unmeasured oxidant
field: low O$_3$ can mark active oxidative chemistry producing toxic secondary products
nobody instruments for. The same argument applies verbatim to PM mass as a proxy for
particle oxidative potential.

\item \textbf{Radon methodology, transferable wholesale.} Giang et al.\ co-located
\textbf{45 consumer radon monitors (15 models $\times$ 3 units)} against a reference for
98 days in an occupied home: 8 of 15 models within $\pm$25\% relative percent error,
worst-case \textbf{$-$47\% to $+$220\%}, intra-model CV mostly under 15\%. Different
analyte, identical QA problem to a consumer PM network --- and a rare published estimate
of \emph{between-model} spread at consumer price points.

\item \textbf{Africa, twice, and both times as a gap statement.} Wright et al.\ find only
13 African case-control studies in a decade (indoor PM$_{2.5}$ 96--177\,$\mu$g/m$^3$,
ambient to 259) and \textbf{none} examining pollution--heat interaction; Hulme et al.\
link home and school exposures for 919 Accra schoolchildren. A third African record
(Lagos speciation, Hopke a co-author) ships on metadata only.

\item \textbf{Nil legs.} Consensus nil for the fifth run; arXiv answered without a 429 but
its newest submission is \textbf{2026-08-20}; OpenAlex 429; trial registry a true zero.
\end{itemize}

\clearpage

% ==================================================================== ANALYTICS
\section*{\sffamily\bfseries\color{Deep}Corpus analytics}
\vspace{-1.5mm}{\color{Amber}\rule{\linewidth}{1.1pt}}\vspace{3mm}

\noindent\includegraphics[width=\linewidth]{\FIGDIR/f1_subtopics.png}
\figcap{Subtopic assignment of the 38 in-scope records, single-label by dominant
contribution. \textit{Exposure assessment \& modelling} (12) and \textit{Sensing,
forecasting \& instrumentation} (11) hold \textbf{61\%} of the issue between them. Only
\textit{Cardiovascular \& metabolic} is vacant; nine or more clusters have been occupied
on ten previous issues, most recently all ten on 27 August. The single \textit{Burden, policy \& mitigation} record
is the China 2035 NDC co-benefit model, and the single \textit{Neuro / mental health}
record is an $n$=44 exploratory PCA, so two of those nine cells are thin.}

\vspace{2mm}
\noindent\includegraphics[width=\linewidth]{\FIGDIR/f2_design_endpoint.png}
\figcap{Left --- study architecture over 23 distinct design labels, 14 of which are new to
the design map and were added to \texttt{plots.py} in the same edit that wrote them.
\textit{Modelling / inventory} (10) and \textit{Measurement campaign} (8) lead;
\textit{Review / synthesis} holds 6 (three systematic reviews, three critical reviews).
Only \textbf{2} records ship on \textit{Metadata only}, against 6 on 2 September.
Right --- \textbf{23 of 38 (61\%)} carry no human health endpoint. That is the largest
\emph{count} the watch has recorded but not the largest share: 20 August reached 72\% at
$n$=18 and 2 August 71\% at $n$=31, and the last three issues have run 58\%, 60\% and
61\%, so the level is a plateau rather than a jump. Read \textit{Cancer} and \textit{Oncologic} as two different
quantities that the canonical map does not merge: the former is 1{,}625 observed incident
cancers in a Canadian birth cohort, the latter a computed carcinogenic-risk index in a
metadata-only Lagos record.}

\vspace{2mm}
\noindent\includegraphics[width=\linewidth]{\FIGDIR/f3_metric_geography.png}
\figcap{Left --- particle metric. \textit{PM (unspeciated) / emissions} at 16 is inflated
by the instrument and chemistry records whose measurand is not a regulated mass fraction
(lidar smoke mass, CCN, INP, organosulfates, filtration efficiency, and the one gas-phase
NO$_2$ sensor carried here). \textit{PM$_{2.5}$ only} at 13 is the honest epidemiological
core. Right --- geography. China 12, North America 8, Europe 5. The
\textit{Global / multi-region} bar of 8 is \textbf{not} eight multi-region studies: four
are genuinely multi-region reviews and four are laboratory records, which this pipeline
routes to that bucket by intent. \textbf{Sub-Saharan Africa at 3 ties the watch record}, set
on 6 August; South Asia (Delhi) and East Asia ex-China (South Korea) contribute one each;
Latin America, the Middle East, Southeast Asia and Oceania are absent.}

\vspace{2mm}
\noindent\includegraphics[width=\linewidth]{\FIGDIR/f6_lifecourse.png}
\figcap{Life-course windows addressed; counts are non-exclusive and cover only the 15
records with a human endpoint. \textbf{Puberty and adolescence is empty} for the second
consecutive issue. Working-age adulthood (4) is the largest cell, carried by two occupational records, the
rural solid-fuel personal-monitoring study and the mTBI cohort; it also led on 29, 30 and
31 August, so this is the resumption of a pattern rather than a new one.}

\vspace{2mm}
\noindent\includegraphics[width=\linewidth]{\FIGDIR/f4_heatmap.png}
\figcap{Subtopic $\times$ study-architecture cross-tabulation. The dense cell is
\textit{Exposure assessment \& modelling} $\times$ \textit{Modelling / inventory}; the
\textit{Sensing} row spreads across measurement campaign, modelling and laboratory
development rather than concentrating, which is what a genuinely instrument-heavy day
looks like as opposed to a day of forecasting papers filed under sensing.}

\vspace{2mm}
\noindent\includegraphics[width=\linewidth]{\FIGDIR/f5_forest.png}
\figcap{Twelve estimates from \textbf{four} sources, and they are not commensurable.
Six are per-IQR component-specific frailty hazard ratios from one CHARLS analysis and are
strongly correlated with one another (the components co-vary in the same TAP exposure
surface). Two are per-10{,}000\,cm$^{-3}$ prenatal-UFP cancer hazard ratios whose
intervals both cross the null. Two are per-IQR ozone rate ratios on 3.48 million ED visits
--- near-null point estimates with intervals so tight they render as bare markers, which
is the correct visual for a 1.9\% excess. The two outliers at the top of the panel are the Guangdong CHD
odds ratios; their intervals are narrower than any biological reading justifies and they
should be treated as a register-scale artefact, not as the day's largest true effect.
Wood-dust spirometry (pooled MD $-$8.57 \%pred FEV$_1$) is absent because a mean
difference has no place on a log-ratio axis.}

\clearpage

% ==================================================================== DIGEST
\section*{\sffamily\bfseries\color{Deep}Paper-level digest}
\vspace{-1.5mm}{\color{Amber}\rule{\linewidth}{1.1pt}}\vspace{1mm}

{\fontsize{7.4}{9}\selectfont\color{Slate}Ordered by cluster. Each entry gives the
methodological core and the finding that matters, not an abstract paraphrase.\par}

\band{Reproductive \& developmental\hfill 2 records}{Sky}

\paperentry{Xu ZL et al.\ 2026}
{Component-specific PM$_{2.5}$ windows and oocyte yield in assisted reproduction}
{Biomed Environ Sci}
{\textbf{51{,}122 ART cycles} with individual PM$_{2.5}$ and its sulfate, nitrate,
ammonium, organic-matter and black-carbon fractions assigned over recent (0--3 months),
distal (4--12 months) and cumulative windows; negative binomial regression on total,
mature and normally fertilised oocyte counts, DLNM for window identification, mixture
models for joint effect. Higher exposure tracked poorer oocyte outcomes throughout, with
\textbf{stronger effects for the distal window} and two sensitive periods --- 1 month and
6--11 months before retrieval. Sulfate and ammonium dominated the mixture. The 6--11 month
window is biologically coherent with primordial-to-antral follicle transit, which is the
paper's strongest argument. Its weakness is that ART patients are a selected population and
the exposure model is residential, so a distal window is also the window over which
residential misclassification is largest.}
{10.3967/bes2026.072}

\paperentry{Chen Y et al.\ 2026}
{PM$_{2.5}$ components in early pregnancy and congenital heart disease}
{Sci Rep}
{\textbf{2{,}601{,}300 live births and stillbirths}, Guangdong 2014--2017, with
satellite-derived component exposures resolved to the week of gestation and quantile
g-computation for mixture weights. Effects peak in \textbf{pregnancy week 1}: sulfate
\textbf{OR 2.23 (2.19--2.27)}, black carbon \textbf{1.95 (1.91--1.98)}, joint weights
86.8\% and 13.2\%. Read the intervals before the point estimates. A near-doubling of CHD
odds with a 95\% CI of width 0.08 implies a precision no exposure model at this spatial
scale supports, and a week-1 peak in a satellite product with month-scale effective
resolution is more consistent with seasonal and spatial confounding than with a
window-specific teratogenic mechanism. The scale of the register is real; the causal
contrast is not established.}
{10.1038/s41598-026-59940-7}

\band{Respiratory \& allergic\hfill 2 records}{Sage}

\paperentry{Wright et al.\ 2026}
{Air pollution and respiratory outcomes in Africa --- a review that mostly documents absence}
{Curr Environ Health Rep}
{PRISMA search of PubMed, ScienceDirect and Elicit for African case-control studies
2015--2025 returned \textbf{13 studies from 10 countries}. Reported concentrations are the
finding: indoor PM$_{2.5}$ \textbf{96--177\,$\mu$g/m$^3$} in biomass-using homes and
ambient PM$_{2.5}$ to \textbf{259\,$\mu$g/m$^3$}, an order of magnitude above WHO guidance.
Nitrogen oxides tracked reduced lung function in children; household air pollution tracked
under-five mortality and low birthweight; associations with acute respiratory infection
were inconsistent across settings. \textbf{Only three studies measured temperature and none
examined heat--pollution interaction}, which is the stated motivation for the review and
which it therefore cannot deliver. Restricting to case-control designs also excludes the
cohort and panel literature, so this is a floor on the African evidence base, not a census.}
{10.1007/s40572-026-00561-7}

\paperentry{Cui L et al.\ 2026}
{Epidemic thunderstorm asthma in China as an \textit{Artemisia} phenotype}
{Clin Rev Allergy Immunol}
{54 primary studies (43 international, 11 Chinese --- Yulin, Hohhot, Chongqing, Lanzhou,
Shenzhen) from 436 records, risk of bias by Newcastle--Ottawa and Murad, certainty by
GRADE. Three patterns recur: Chinese events are associated with \textbf{mugwort
(\textit{Artemisia})} rather than ryegrass sensitisation; \textbf{40--60\% of Chinese
cases and up to 78\% of Australian and UK cases had no prior asthma diagnosis}; and
individual Chinese localities report recurrence. No meta-analysis, and GRADE certainty for
direct pollen rupture, pollution-mediated severity and pre-emptive treatment is
\textbf{low to very low}. The measurement gap is explicit and is the actionable part:
Chinese reports identify sensitisation only to genus level and \textbf{do not speciate the
airborne pollen}, so no aerobiological trigger threshold can be set. The authors' first
research priority is field aerosol sampling.}
{10.1007/s12016-026-09194-y}

\band{Mechanistic toxicology\hfill 3 records}{Violet}

\paperentry{Na HW et al.\ 2026}
{Airborne metals reconstructed from worker sebum, then dosed onto skin}
{J Toxicol}
{ICP-MS on sebum from \textbf{62 indoor and 40 outdoor workers} in Shanghai across 24
metals, with large outdoor--indoor contrasts (cobalt 925$\times$, arsenic 135$\times$,
barium 29$\times$, manganese 17$\times$). Sebum-mimicking mixtures were then formulated at
the measured profiles and applied to keratinocytes, sebocytes and ex vivo scalp skin: the
outdoor mixture reduced viability, raised lipid peroxidation, induced $\gamma$H2AX foci,
upregulated \textit{IL1A}, \textit{IL1B}, \textit{TNFA}, \textit{IL6} and \textit{PTGS2},
and activated NF-$\kappa$B and MAPK. The design is the interesting part --- a biomonitoring
matrix used to specify the in vitro dose rather than an arbitrary slurry concentration.
The weakness is that sebum metal load is a cumulative deposition proxy of unknown
relation to inhaled or dermally absorbed dose, and indoor/outdoor is a crude exposure
contrast with occupation confounded into it.}
{10.1155/jt/7925949}

\paperentry{Wang Y et al.\ 2026}
{$\alpha$-Dicarbonyl aqueous products modulate cellular ROS in opposite directions}
{Environ Sci Technol Lett}
{Glyoxal or methylglyoxal reacted with ammonium sulfate or glycine in the aqueous phase;
HRMS shows ammonium routes yield imidazole-type nitrogen heterocycles and glycine routes
yield more pyrrole-related species. Alone, neither product set induces much ROS in A549
cells. Co-exposed with 1,4-naphthoquinone, Cu$^{2+}$ or both, \textbf{imidazole-dominated
mixtures enhance ROS and pyrrole-rich mixtures suppress it}, with acellular assays
attributing the split to promotion versus suppression of quinone redox cycling and
metal-associated pathways. This is a direct challenge to any oxidative-potential assay
read as a scalar: the same secondary organic aerosol aging pathway can raise or lower
measured ROS depending on the nitrogen source. Cell-line work at unstated
atmospheric relevance of concentration, so treat as mechanism, not as an exposure-response.}
{10.1021/acs.estlett.6c00722}

\paperentry{Cao et al.\ 2026}
{Black carbon interfacial chemistry and multiscale aging}
{Environ Sci Atmos}
{A review that puts BC \emph{particle state} --- surface reactivity, coating thickness,
morphology, aging timescale --- rather than BC mass at the centre, and follows it through
to optics, cloud activation, removal, lifetime and toxicity. The framing matters for
instrumentation because every one of those downstream properties is a function of a
variable that filter-based absorption and single-particle soot photometry constrain only
partially, and that no low-cost optical sensor constrains at all. No primary data and no
quantitative synthesis; the Crossref deposit carries only a graphical-abstract blurb, so
the parameterisation recommendations were not verifiable at digest time. Listed as a
reference framework rather than as evidence.}
{10.1039/d6ea00093b}

\band{Occupational \& indoor\hfill 3 records}{Clay}

\paperentry{Srinivasan et al.\ 2026}
{Occupational wood dust and lung function --- pooled, and heterogeneous}
{Natl Med J India}
{PRISMA review of PubMed and Google Scholar; 55 studies qualitatively, 16/14/9 pooled by
outcome. Random-effects pooled mean differences in percent predicted:
\textbf{FEV$_1$ $-$8.57 (95\% CI $-$12.07 to $-$5.08)}, \textbf{FVC $-$7.43
($-$11.35 to $-$3.51)}, \textbf{PEFR $-$13.43 ($-$25.14 to $-$1.72)}. Sawmill workers
trended obstructive, furniture workers restrictive. Three things limit it:
\textbf{I$^2$ exceeds 95\%} on every pooled outcome, the funnel plot indicates publication
bias, and exposure is binary --- no study contributed a dust concentration, so there is no
exposure--response and no basis for an occupational limit. Searching only PubMed and
Google Scholar is thin for an occupational literature that sits substantially in
non-indexed regional journals.}
{10.25259/nmji\_668\_2024}

\paperentry{Feng R et al.\ 2026}
{Clean heating renovation cuts nitro-PAHs but not nitrated phenols}
{Environ Sci Technol Lett}
{Size-resolved \emph{personal} exposure to particulate nitrated phenols (NPs) and
nitro-PAHs among rural users of clean coal, raw coal chunk and biomass, paired with A549
cytotoxicity and NF-$\kappa$B reporter assays. Biomass users carried up to
\textbf{11.5$\times$} the PM$_{2.5}$-bound n-PAH exposure of clean-coal users, while NPs
varied \textbf{$\leq$1.3-fold} across fuels (clean coal 64.0$\pm$12.1 vs raw coal
57.4$\pm$25.8\,ng/m$^3$, $p$=0.545). Toxicologically the two classes diverge as well:
n-PAHs were more cytotoxic by IC$_{50}$, NPs activated NF-$\kappa$B more strongly via
p-I$\kappa$B-$\alpha$ and COX-2. The policy reading is sharp --- a fuel switch evaluated on
PM mass or on PAHs alone will score a benefit that the NP fraction does not deliver.
Group sizes are not stated in the abstract and the in vitro work uses representative
compounds, not the collected particles.}
{10.1021/acs.estlett.6c00629}

\paperentry{Jamali et al.\ 2026}
{All-cellulose HEPA nanofilter in a wooden frame}
{Int J Biol Macromol}
{Fog-sprayed, freeze-dried bacterial-cellulose nanofibres on a wood-derived
cellulose-microfibre tissue substrate, in a modular wooden frame with a starch-adhesive
gasket. Removal efficiency rises with BCNF grammage to \textbf{2\,g/m$^2$} and falls above
it --- consistent with pore closure raising velocity through the remaining paths. Optimum
gives \textbf{99.998\% PM$_{0.3}$ at 150\,Pa}, inside the EN~1822 H13--H14 band, and
PM$_{0.3}$ capture \emph{peaks at high RH} under cyclic humidity, the opposite of the
usual electret degradation. The claim being made is biodegradability at H13 performance,
and the tests reported are single-pass efficiency and pressure drop; there is no loading
lifetime, no microbial growth assessment on a cellulose substrate held at high RH, and no
mechanical durability data. Promising material, incomplete filter.}
{10.1016/j.ijbiomac.2026.154178}

\band{Sensing, forecasting \& instrumentation\hfill 11 records}{Amber}

\paperentry{Kaur et al.\ 2026}
{Plantower PMSX003N against FEM through a Utah winter}
{Aerosol Sci Technol}
{Three PMSX003N, three PMS5003T and three OPC-N3 co-located with FEM monitors at a Utah
DEQ station for five months, plus laboratory size-binning against monodisperse DOS
aerosol. The laboratory result drives everything: the PMSX003N
\textbf{preferentially assigns particles $>$1\,$\mu$m to the 2.5--5\,$\mu$m bin}. In dust
events this recovers coarse mass its predecessor loses --- \textbf{PM$_{2.5}$ slope
1.85--1.91, $R^2>0.7$} vs \textbf{0.15, $R^2<0.37$}; \textbf{PM$_{10}$ slope 0.91--1.07}
vs \textbf{$\sim$0.02} --- and under freezing fog, mist, snow and rain it substantially
overestimates both fractions. Three units per model over one site-season is a small basis
for a between-unit claim, and a slope near 1.9 is a large residual bias even where the
correlation is good. Humidity-dependent correction is not optional for this part.}
{10.1080/02786826.2026.2716649}

\paperentry{Guth et al.\ 2026}
{Time-resolved emission factors for 19 prescribed burns, measured with low-cost sensors}
{Atmosphere}
{Emission factors for CO, PM$_{2.5}$, EC, OC and speciated organics across \textbf{19
prescribed fires} in Colorado and southeastern Georgia, using multiple low-cost monitors
per burn to get temporal resolution that filter-based EF campaigns cannot reach. CO EFs
are driven by combustion phase --- \textbf{+3260\%} at low modified combustion efficiency,
i.e.\ smouldering. Broadcast and pile burns are compared directly, which is rare. The
methodological point for a distributed network is that the LCS bias that disqualifies a
sensor for regulatory-equivalent concentration is tolerable here because the measurand is a
ratio within a plume; the risk is that any RH- or size-dependent bias in the PM channel
propagates straight into the PM$_{2.5}$/CO EF ratio, and the abstract does not report the
sensor correction applied.}
{10.3390/atmos17090867}

\paperentry{Xing et al.\ 2026}
{DeepAQM --- observation-constrained ML surrogate for the NOAA forecast system}
{Environ Sci Technol}
{A ConvLSTM--ResNet fusion reconstructs spatially continuous surface PM$_{2.5}$ and O$_3$
from VIIRS AOD, TEMPO NO$_2$ and EPA AirNow, and supplies them as initial conditions for
72\,h UFS-AQM forecasts over the continental US. It reproduces UFS-AQM spatial patterns
($R^2>0.9$ for O$_3$, to 0.80 for PM$_{2.5}$, normalised mean bias within $\pm$10\%).
The result that matters: observation-updated initial conditions lift site-level $R^2$ in
the first 10 forecast hours from \textbf{$\sim$0.10--0.15 to 0.40--0.50 for PM$_{2.5}$},
and observation-based fine-tuning carries overall $R^2$ to \textbf{0.41 for PM$_{2.5}$}
across the full 72\,h. Note what that number is: a fourfold gain at short lead, decaying
to a forecast that still explains under half the variance. TEMPO makes this a
current-generation result rather than a reanalysis exercise.}
{10.1021/acs.est.6c04600}

\paperentry{Veselovskii et al.\ 2026}
{Five-channel fluorescence lidar retrieves smoke mass inside the boundary layer}
{Atmos Meas Tech}
{Smoke and urban aerosol have distinct fluorescence spectra, and a multi-channel
fluorescence lidar exploits that to separate them and retrieve smoke mass concentration
with an estimated detection threshold of \textbf{$\sim$0.1\,$\mu$g/m$^3$}. Over Moscow,
smoke was detected in \textbf{59 of 67 sessions between April and October 2024}, mostly
from long-range Atlantic transport; only 12 episodes traced to southern Russian fires.
Within the PBL, transported North American wildfire smoke contributed
\textbf{$\sim$1\,$\mu$g/m$^3$} against \textbf{up to 50\,$\mu$g/m$^3$} for regional fires.
This is the source-resolved column-to-surface capability that mass-only networks lack, and
0.1\,$\mu$g/m$^3$ is a genuinely low threshold. Single site, single instrument, and mass
retrieval requires an assumed fluorescence-to-mass conversion whose uncertainty is the
main open term.}
{10.5194/amt-19-5587-2026}

\paperentry{Zhang Z et al.\ 2026}
{Microfluidic multispectral diffraction for black carbon quantification}
{Anal Chem}
{A microfluidic chip does inertial separation, low-velocity sheath-flow regulation and
high-voltage electrostatic capture to enrich submicron BC in the imaging region
(\textbf{92\% enrichment efficiency}); a lensless diffraction system then acquires
hyperspectral fingerprints over 400--800\,nm, which a graph model maps to concentration.
Over \textbf{0.01--0.20\,mg/m$^3$}, $R^2$=0.98. The discriminating claim is more
interesting than the calibration: the spectral fingerprint separates BC from
similarly-sized microplastics and pathogenic spores, which absorption-based
instruments cannot do. Bench work only --- no ambient co-location, no interference testing
against real coated or aged BC, and 0.01\,mg/m$^3$ (10\,$\mu$g/m$^3$) is a high floor for
ambient BC, which typically runs 0.5--5\,$\mu$g/m$^3$.}
{10.1021/acs.analchem.6c03051}

\paperentry{Giang et al.\ 2026}
{Fifteen consumer radon monitors, 98 days, one occupied house}
{J Radiol Prot}
{45 units (15 models $\times$ 3) co-located with an AlphaGuard PQ2000 and two alpha-track
detectors in the lowest lived-in space of an occupied Ontario home, with PM$_{2.5}$,
CO$_2$, temperature and RH logged alongside; reference mean \textbf{99\,Bq/m$^3$} (range
62--141). Cumulative performance: \textbf{8 of 15 models} had two or more units within
$\pm$25\% relative percent error, the best reaching 3--5\% after 13 weeks, while others
ran from \textbf{$-$47\% to $+$220\%}. Intra-model CV was mostly under 15\%; real-time
RMSE for the best units was 5--6\,Bq/m$^3$ with Spearman $\rho$=0.70--0.77 on two-day
averages. Different analyte, and the QA design is directly transferable: this is the
between-model spread estimate that consumer PM networks routinely lack. One home, one
season, one concentration regime.}
{10.1088/1361-6498/aea16a}

\paperentry{Cui C et al.\ 2026}
{Eighteen years of CCN reconstructed by XGBoost with SHAP attribution}
{Environ Sci Technol}
{An observation-driven XGBoost framework reconstructs CCN number concentration at a North
China Plain mountain site, 2007--2025. Particle number size distribution carries
\textbf{55.7--66.3\%} of predictive importance, with the controlling size range shifting
from 100--200\,nm at 0.2--0.4\% supersaturation to \textbf{50--100\,nm at 1.0\%}. No
monotonic CCN trend, but summer bulk activation ratio and hygroscopicity parameter both
\textbf{declined from 2009 to 2025}. Coupling PMF with the ML model attributes that to
falling SO$_2$ influence and rising contributions from new particle formation and growth
--- a legible consequence of desulfurisation policy. Reconstruction over an 18-year record
inherits every instrument change in that record, and SHAP importance is not a causal
decomposition.}
{10.1021/acs.est.6c04501}

\paperentry{Galioto et al.\ 2026}
{Physics-informed ML for street-canyon NO$_2$ in Palermo}
{Atmosphere}
{Six years of hourly data (2020--2025) on one of Europe's densest traffic corridors,
resolved to \textbf{$\sim$11\,m}: COPERT-Italy emission factors over the OpenStreetMap
network, dispersed with the SIRANE street-network model, then fused with measurements in
an ML stage and mapped in GIS. Observation--model scatter over \textbf{44{,}752 hours}
centres on the 1:1 line. The genuinely useful move is inverting the chain to recover
corridor traffic volume from observed NO$_2$ --- a template for using a fixed monitor as a
traffic sensor, which transfers to PM only where the non-exhaust fraction is separately
handled. Analyte is gas-phase; single corridor; and a physics-informed hybrid validated
against the same stations that trained it needs an out-of-corridor test before the
inversion can be trusted.}
{10.3390/atmos17090865}

\paperentry{Shafi \& Scafetta 2026}
{Reconstructing a decade of missing Delhi monitoring records}
{Environ Sci Pollut Res Int}
{Delhi went from \textbf{4 stations in 2014 to 45 in 2024}, leaving a fragmented record
that defeats trend analysis. For each missing daily value across six pollutants, the
workflow selects the four most-correlated stations as predictors, evaluates \textbf{35
regression models} in MATLAB's Regression Learner, takes the best, and iterates until the
matrix is full. Tree ensembles and Gaussian/rational-quadratic kernels consistently beat
multilinear regression; artificial-gap experiments validate across pollutants, stations
and years, best for PM$_{2.5}$ and other regionally-driven species. The structural problem
is unaddressed and important: stations are missing non-randomly, and imputing from
correlated neighbours smooths exactly the local excursions that a network is deployed to
detect. Reconstructed series should not be used for exceedance statistics.}
{10.1007/s11356-026-38181-1}

\paperentry{Alderotti et al.\ 2026}
{Raman spectroscopy as a non-destructive biomonitoring readout}
{Sci Total Environ}
{Mature \textit{Quercus ilex} sampled along a high/intermediate/low NO$_2$ and PM$_{10}$
gradient in Tuscany. Leaf Raman spectra show coordinated modulation of pigment
(chlorophyll, carotenoid) and inducible flavonoid bands, cross-validated against
destructive biochemistry and chlorophyll fluorescence, and consistent with
\emph{acclimation} rather than photoinhibitory damage. A composite Raman index tracked
site-level NO$_2$ most closely, which the authors read as nitrogen driving the response.
The attraction is a rapid, non-destructive, scalable field readout; the limitation is that
a three-site gradient cannot separate NO$_2$ from PM$_{10}$ from any co-varying urban
factor, so the index is at best a composite urban-exposure indicator and not a
pollutant-specific one.}
{10.1016/j.scitotenv.2026.182253}

\paperentry{He J et al.\ 2026}
{Room-temperature ppb NO$_2$ chemiresistor by molecular packing control}
{ACS Sens}
{Alkyl chains of varying length on cyano-substituted distyrylbenzene tune molecular
conformation and film microcrystallinity, exposing more surface-accessible functional
groups and raising NO$_2$ sensitivity; the optimised device reaches a \textbf{30\,ppb}
detection limit at room temperature, and response is \emph{enhanced} by humidity through
NO$_2$--H$_2$O synergy. Demonstrated in a flexible wireless vehicle-exhaust monitor.
Carried here on the sensing axis with the analyte stated plainly: this measures a gas, not
a particle, and contributes nothing to PM mass or size. It is included because
room-temperature operation and positive rather than negative humidity dependence are both
unusual for the co-located gas channels that sit in multi-pollutant low-cost nodes. No
long-term drift or cross-sensitivity data.}
{10.1021/acssensors.6c00977}

\band{Exposure assessment \& modelling\hfill 12 records}{Amber}

\paperentry{Bravo et al.\ 2026}
{Chicago's 2023 wildfire smoke, by tract}
{J Expo Sci Environ Epidemiol}
{Hourly PM$_{2.5}$ on a \textbf{500\,m} grid aggregated to daily 24-hour means for all
\textbf{792 Chicago census tracts} in nine city regions, 1 May--31 July 2023, examined
against demographic, socio-economic and health profiles. The title carries the finding:
during the long-range wildfire-smoke episodes, the usual within-city exposure gradient
compressed --- smoke reached every region --- while the pre-existing disparity in baseline
PM$_{2.5}$ and in underlying health persisted underneath it. That combination is the
argument for source-resolved rather than total-mass exposure metrics in environmental
justice work, since a total-mass summer average would show convergence and hide the
differential burden. Modelled rather than measured exposure at tract level, one summer, one
city, and no health outcome is linked here.}
{10.1038/s41370-026-00956-6}

\paperentry{Hulme et al.\ 2026}
{Home and school exposure, jointly, for 919 Accra schoolchildren}
{Environ Res Lett}
{Land-use regression annual means for PM$_{2.5}$, NO$_2$ and BC plus L$_\mathrm{day}$,
L$_\mathrm{night}$ and L$_\mathrm{den}$ noise, linked to geocoded home and school points
for \textbf{919 of 1{,}034} children aged 8--12 across 90 Accra schools, then combined into
a time-location-weighted total. Median home PM$_{2.5}$ \textbf{33.8\,$\mu$g/m$^3$},
NO$_2$ 55.5, BC 5.2$\times$10$^{-5}$\,m$^{-1}$; school values slightly lower
(29.9 / 53.5 / 4.9). Time-weighted totals exceeded WHO guidance for most children, and
noise exceeded WHO road-traffic thresholds at home (L$_\mathrm{den}$ 66.3) and Ghana's
educational-area standard at school. The joint air-plus-noise, home-plus-school design is
the contribution --- both are under-measured in sub-Saharan Africa. Annual LUR means
cannot resolve the school day, so the time-weighting is coarser than it looks.}
{10.1088/1748-9326/ae8d40}

\paperentry{Liu S et al.\ 2026}
{Wind-shear regimes and the vertical PM$_{2.5}$ profile in a basin}
{Environ Pollut}
{Six instrument types --- surface network, wind profiler radar, microwave radiometer,
aerosol lidar, radiosonde and UAV --- plus WRF, over Taiyuan Basin. Unsupervised clustering
gives three vertical-wind-shear regimes: strong directional near-surface shear with
near-surface accumulation and the shallowest boundary layer; speed shear above 500\,m with
an \emph{elevated} PM$_{2.5}$ maximum near that height; and deep shear below 2000\,m with
stronger winds and low concentrations. UAV profiles over 0--500\,m confirm the three
shapes. The 17 January 2024 episode makes the operational point: strong shear did not
enhance vertical exchange while the shear layer sat above a shallow stable layer, and only
coupled once daytime heating deepened the boundary layer into it. Directly relevant to
siting --- a surface monitor cannot see regime 2 at all. One basin, one winter.}
{10.1016/j.envpol.2026.129103}

\paperentry{Yu H et al.\ 2026}
{Sea--land breeze recirculation of black carbon in Qingdao}
{Environ Sci Technol}
{A full year of online BC (2021) split into fossil-fuel and biomass-burning fractions by
the Aethalometer model, then attributed to meteorology by random forest and trajectory
analysis. Annual mean \textbf{1.78$\pm$1.29\,$\mu$g/m$^3$}; fossil fuel \textbf{88.2\%} of
mass, biomass burning rising to \textbf{19\% in winter}. The two fractions respond
differently --- BC$_\mathrm{ff}$ to wind-driven ventilation and transport, BC$_\mathrm{bb}$
to low temperature --- and BC enhancement under sea--land breeze peaks at
\textbf{2.0--3.5\,m/s}, with diurnal persistence and directional shifts consistent with
return transport and reaccumulation. Coastal siting guidance follows directly: a single
coastal monitor under SLB conditions is sampling a recirculated air mass, not a fresh one.
Aethalometer-model source splitting depends on assumed absorption {\AA}ngstr\"om
exponents, which is the dominant uncertainty here.}
{10.1021/acs.est.6c10060}

\paperentry{Wu Q et al.\ 2026}
{Nighttime new particle formation, and it is not always sulfuric acid}
{Environ Sci Technol}
{Roughly \textbf{1.5 years} of molecular-level measurement in coastal Dalian. Spring and
winter nocturnal NPF is driven by H$_2$SO$_4$, whose dimer formation proceeds near the
collision limit. In summer and warm autumn, H$_2$SO$_4$ and its dimers are scarce while
oxygenated organic molecules are abundant, and \textbf{OOM dimers correlate with nascent
particles at mean $r$=0.55}. Organic-driven nucleation has been reported for biogenic
environments; showing it under anthropogenic influence at night is the novel claim.
Implication for number-concentration monitoring: an ultrafine source term exists at night
that photochemically-anchored parameterisations will miss. Correlation of 0.55 is
suggestive rather than mechanistic, and one coastal site cannot establish how general the
seasonal switch is.}
{10.1021/acs.est.6c04493}

\paperentry{Gao J et al.\ 2026}
{Are inverse ozone--health associations an artefact of measuring the wrong oxidant?}
{Environ Sci Technol}
{A perspective, not a study. Population studies not infrequently find \emph{lower} O$_3$
associated with \emph{higher} disease risk. NO titration near traffic and differential
measurement error are the standard explanations and do not fully account for it. The
alternative advanced here: low measured O$_3$ can mark locations and times of active
oxidative chemistry --- O$_3$ being consumed, or parallel oxidant pathways running ---
which generate toxic secondary products that no monitoring network measures. If so, O$_3$
is functioning as a negatively-signed proxy for an unmeasured oxidant field. The same
logic transfers directly to PM mass as a proxy for oxidative potential, which is why this
is carried in a PM watch. No data and no test proposed beyond a call for expanded
measurement, so it stands or falls on plausibility.}
{10.1021/acs.est.5c17322}

\paperentry{Yang \& Li 2026}
{Modelling the whole indoor-radon distribution, not the area mean}
{Environ Sci Technol}
{Over \textbf{6 million} short-term radon measurements across the contiguous US,
2000--2021, used to predict \emph{percentiles} at ZIP Code Tabulation Area level and then
to model the within-ZCTA distribution rather than summarising each area by a mean or
median. The consequence is quantitative: roughly \textbf{half of US ZCTAs have at least one
in five households} above policy-relevant thresholds --- a figure an area-average surface
cannot produce, because the exceeding fraction is a property of the spread. Same argument
applies to any PM exposure surface used for threshold-crossing statistics: modelling the
mean and modelling the tail are different tasks. Radon test data are self-selected
(homeowners who chose to test, often at point of sale), which biases both the level and
the spread in unknown directions.}
{10.1021/acs.est.6c01333}

\paperentry{Lai et al.\ 2026}
{Two-thirds of ambient organosulfates have no assigned source}
{Environ Sci Technol}
{Organosulfates contribute up to \textbf{30\% of aerosol mass}. Of \textbf{366} ambient OS
molecular formulas compiled from field studies, only about \textbf{36\%} map to a known
precursor or formation mechanism, leaving \textbf{64\% unclassified}. The proposal here is
that some unclassified species are not freshly formed at all but higher-generation products
of heterogeneous and aqueous multiphase transformation of existing OS: identified
transformation products already account for \textbf{$\sim$17\%} of reported formulas and
explain an additional \textbf{5\%} of the unclassified set. Modest arithmetic, but it
reframes a source-apportionment problem as an aging problem. Compilation-based, dependent
on the sampling and HRMS coverage of the underlying field studies, and 5\% is a small
demonstration of a mechanism claimed to be general.}
{10.1021/acs.est.6c08394}

\paperentry{Schaefer et al.\ 2026}
{Deep convection sorts ice-nucleating particles by activation temperature}
{Atmos Chem Phys}
{Immersion-mode INP filter samples from the HALO aircraft up to \textbf{14.5\,km} during
CIRRUS-HL (summer 2021), taken in the inflow and outflow of deep convective clouds and
from cloud-particle residuals. Inflow spectra contain INPs active at both high
($T>-15^{\circ}$C) and low ($T<-20^{\circ}$C) temperatures; \textbf{outflow spectra contain
only the low-temperature population}. The reading is precipitation scavenging in the
updraft removing the warm-active INPs while the cold-active ones are carried efficiently
into the free troposphere. Cloud microphysics rather than exposure science, carried for the
particle-transport result: convection is a size- and composition-selective vertical filter,
which matters for any attempt to relate column-integrated aerosol to surface composition.
Offline analysis, filter sampling, one campaign.}
{10.5194/acp-26-12505-2026}

\paperentry{Wang X et al.\ 2026}
{Coupling coordination between bus-stop dust and soil metals}
{Environ Monit Assess}
{52 paired surface-dust and soil samples (104 total) at bus stops in Weiyang District,
Xi'an, assayed for Cr, Mn, Cu, Zn, As and Pb and scored with Igeo and the Nemerow index.
All but As and Mn exceeded Shaanxi soil background. Spatial patterns differ by medium ---
dust higher east, soil higher south --- and sources resolve to traffic, natural and
industrial. Hand-to-mouth ingestion dominated the exposure route and non-carcinogenic risk
was higher for children than adults. The novel element is applying a coupling-coordination
degree between the two media, which came out \emph{moderate} and is read as partial
interactive migration with shared sources. That index is a descriptive statistic borrowed
from regional economics, not a transport model, so ``interactive migration'' is an
interpretation rather than a demonstrated flux.}
{10.1007/s10661-026-15820-1}

\paperentry{Pach\'es et al.\ 2026}
{Regional background PM$_{2.5}$ in the Valencian Community by clustering}
{Atmos Pollut Res}
{Ships on metadata only. The Elsevier deposit carries no abstract, the record is not yet in
PubMed, Europe PMC or Semantic Scholar, and ScienceDirect refused retrieval from this
sandbox, so the network size, the clustering method and the resulting background estimates
are all unavailable. Counted and listed; no numbers claimed. Flagged for the abstract retry
list. Background/local-increment separation is directly relevant to siting a low-cost
network --- the increment is what a dense network is deployed to resolve --- so this is
worth revisiting once the abstract deposits.}
{10.1016/j.apr.2026.103192}

\paperentry{Odu-Onikosi et al.\ 2026}
{PM$_{2.5}$-bound toxic elements in Lagos}
{Atmos Pollut Res}
{Ships on metadata only, same Elsevier deposit gap as above: no abstract, not in
PubMed or Europe PMC, ScienceDirect blocked. Sampling duration, element list, apportionment
method and risk estimates are all unretrievable at digest time. Counted and listed; no
numbers claimed, and flagged for the abstract retry list. Noted because the author list
includes P.\,A.\ Solomon and P.\,K.\ Hopke, and because speciated PM$_{2.5}$ from Lagos ---
a megacity of roughly 15 million with almost no published speciation record --- is exactly
the geographic gap the two African records above document.}
{10.1016/j.apr.2026.103193}

\band{Other clinical endpoints\hfill 3 records}{Slate}

\paperentry{Batisse et al.\ 2026}
{Prenatal ultrafine particles and childhood cancer in 965{,}000 children}
{Environ Epidemiol}
{Cohort of over \textbf{965{,}000} children born in Montreal or Toronto 1999--2020, with
high-resolution outdoor UFP number concentration \emph{and mean UFP size} assigned to
six-digit residential postal codes from annual tax records, followed to age 14 (mean 9.1
years) against the Canadian Cancer Registry. \textbf{1{,}625} incident cancers. Per
10{,}000\,particles/cm$^3$ of prenatal exposure: acute lymphoblastic leukaemia
\textbf{HR 1.12 (0.94--1.34)}, neuroblastoma \textbf{1.15 (0.85--1.57)}; both intervals
include the null, both persisted across exposure windows and sensitivity analyses, and
other sites were null. Adjusting for UFP size as well as number is the methodological
advance --- number and size are correlated and separately toxicologically relevant. The
honest summary is a consistent, imprecise positive signal, not a demonstrated association.}
{10.1097/ee9.0000000000000528}

\paperentry{Chen K et al.\ 2026}
{PM$_{2.5}$ components and incident frailty in CHARLS}
{Sichuan Da Xue Xue Bao Yi Xue Ban}
{CHARLS 2011--2020, \textbf{$n$=9{,}603} middle-aged and older adults, frailty index
outcome, exposures from the Tracking Air Pollution product, time-dependent Cox models. Per
IQR of the preceding year: PM$_{2.5}$ \textbf{HR 1.100 (1.015--1.193)}; by component,
organic matter \textbf{1.142 (1.055--1.236)}, black carbon \textbf{1.131 (1.041--1.229)},
ammonium \textbf{1.098 (1.010--1.194)}, sulfate 1.088 (0.996--1.187) and nitrate 1.081
(0.996--1.172). Carbonaceous components rank above bulk mass and above the secondary
inorganic ions, and the two that cross the null are exactly the two secondary ions ---
an ordering consistent with combustion-source toxicity rather than with mass. The
components co-vary strongly in the same exposure surface, so the ranking is suggestive and
the six estimates are not six independent tests.}
{10.12182/20260760103}

\paperentry{Park J et al.\ 2026}
{Ozone, disability and emergency department use in Korea}
{Int J Epidemiol}
{Time-stratified case-crossover on \textbf{3{,}484{,}877} ED visits from the Korean
National Health Insurance Service (2015--19, 2022--23), exposure the 3-day moving average
of daily 8-hour maximum O$_3$, adjusted for daily mean temperature and holidays. Per IQR
($\sim$22\,ppb): \textbf{+1.94\% (1.64--2.25)} in persons with disabilities during the warm
season and +1.93\% (1.62--2.25) year-round, against \textbf{+1.36\% (0.92--1.80)} in
persons without. A roughly 40\% larger relative effect in a population defined
administratively rather than clinically, with an interaction term tested across temperature
strata. No particulate metric is included, which is both the limitation and, read against
Gao et al.\ in this issue, part of the point: a disability register can identify a
susceptible group that pollutant-specific monitoring networks were never designed to
serve.}
{10.1093/ije/dyag189}

\band{Neuro / mental health\hfill 1 record}{Deep}

\paperentry{Paulino et al.\ 2026}
{Exposome burden and cognitive recovery after mild TBI}
{Int J Neurosci}
{\textbf{44} adults with mild TBI (GCS 13--15, 50\% female) prospectively enrolled at an
academic clinic; 15 a priori exposome variables reduced by forward selection to 13 and
summarised by PCA into a cumulative burden score, against residualised 6-month MoCA
dichotomised into worse- and not-worse-than-expected recovery. The latent axis explains
\textbf{16.0\%} of variance and is loaded by education, area deprivation, insurance
coverage, \textbf{PM$_{2.5}$} and adverse childhood events; burden was higher in the
worse-than-expected group (median PC1 1.34 vs $-$0.76, $p$=0.010, holding under
permutation). Carried for the PM$_{2.5}$ loading only. At $n$=44 with 13 variables and an
outcome-informed selection step, this is hypothesis-generating; the loading structure is
mostly socioeconomic and PM$_{2.5}$ may be riding on it.}
{10.1080/00207454.2026.2728613}

\band{Burden, policy \& mitigation\hfill 1 record}{Coral}

\paperentry{Liu Z et al.\ 2026}
{Health co-benefits of China's 2035 NDC arrive before the climate benefits do}
{Environ Sci Technol}
{Integrated assessment over 2025--2060 puts the NDC's net benefit at \textbf{+1{,}500.5
billion USD}, decomposing into \textbf{3{,}924.3 billion} in long-run climate benefits and
\textbf{1{,}659.5 billion} in health co-benefits. The finding is the timing and the
distribution, not the total: in the near term the NDC \emph{reduces} climate benefit,
especially in underdeveloped regions, because coal generation balances renewable
intermittency and rapid afforestation releases soil organic carbon; long-run benefits reach
all regions but stay concentrated in high-income ones. Air-quality health co-benefits are
the component that accrues locally and immediately, which is the argument for treating PM
monitoring as mitigation-policy infrastructure. Monetised co-benefits inherit every
assumption in the concentration--response and valuation chain and should be read as
scenario contrast, not as a forecast.}
{10.1021/acs.est.6c06024}

\clearpage

% ==================================================================== TABLE
\section*{\sffamily\bfseries\color{Deep}Corpus register}
\vspace{-1.5mm}{\color{Amber}\rule{\linewidth}{1.1pt}}\vspace{2.5mm}

{\fontsize{6.7}{8.4}\selectfont
\setlength{\tabcolsep}{2.4pt}
\begin{longtable}{@{}>{\RaggedRight}p{25mm}>{\RaggedRight}p{41mm}>{\RaggedRight}p{27mm}%
>{\RaggedRight}p{23mm}>{\RaggedRight}p{18mm}>{\RaggedRight}p{15mm}@{}}
\toprule
\textsf{\textbf{First author}} & \textsf{\textbf{Focus}} & \textsf{\textbf{Journal}} &
\textsf{\textbf{Design}} & \textsf{\textbf{Metric}} & \textsf{\textbf{Region}}\\
\midrule
\endfirsthead
\toprule
\textsf{\textbf{First author}} & \textsf{\textbf{Focus}} & \textsf{\textbf{Journal}} &
\textsf{\textbf{Design}} & \textsf{\textbf{Metric}} & \textsf{\textbf{Region}}\\
\midrule
\endhead
\multicolumn{6}{@{}l@{}}{\cellcolor{Mist}\textsf{\textbf{\textcolor{Deep}{Neuro / mental health}}}}\\[0.6mm]
\href{https://doi.org/10.1001/jamanetworkopen.2026.25015}{Dales} & Wildfire-sourced PM$_{2.5}$ and ED visits for migraine and other primary headache & \textit{JAMA Netw Open} & Case-crossover & PM$_{2.5}$ (wildfire) & Canada\\
\href{https://doi.org/10.1093/ageing/afag215}{Glover} & Umbrella meta-review of urban environmental exposures and cognitive health & \textit{Age Ageing} & Umbrella review / meta-analysis & PM$_{2.5}$ & Global\\
\href{https://doi.org/10.1016/j.nsa.2026.107019}{Goossens} & In utero air pollution exposure and pre-/postnatal brain development & \textit{Neurosci Appl} & Narrative review & PM$_{2.5}$ / PM$_{10}$ / PM$_{\mathrm{coarse}}$ & Global\\
\href{https://doi.org/10.1016/j.isci.2026.116617}{Ma} & PM$_{2.5}$-induced neurotoxicity and depression: exposure risks and mechanisms & \textit{iScience} & Narrative review & PM$_{2.5}$ + components & Global\\
\href{https://doi.org/10.3760/cma.j.cn112150-20250804-00752}{Pan} & Maternal PM exposure and offspring neurodevelopment: effects and mechanisms & \textit{Zhonghua Yu Fang Yi Xue Za Zhi} & Narrative review & PM (bulk) & China\\
\addlinespace[0.9mm]
\multicolumn{6}{@{}l@{}}{\cellcolor{Mist}\textsf{\textbf{\textcolor{Teal}{Cardiovascular \& metabolic}}}}\\[0.6mm]
\href{https://doi.org/10.3760/cma.j.cn112338-20251120-00843}{Huang} & PM$_{2.5}$ and dyslipidaemia in rural Chinese children; dietary effect modification & \textit{Zhonghua Liu Xing Bing Xue Za Zhi} & Cross-sectional (multistage) & PM$_{2.5}$ & China\\
\href{https://doi.org/10.32394/pe/218966}{Kazmierska} & Air pollution and cardiovascular disease: review of current data & \textit{Przegl Epidemiol} & Narrative review & PM$_{2.5}$ / PM$_{10}$ & Global\\
\href{https://doi.org/10.1007/s00125-026-06814-2}{Xi} & Air pollution, built environment and temperature vs. prediabetes phenotypes & \textit{Diabetologia} & Prospective cohort & PM$_{2.5}$ & Germany\\
\href{https://doi.org/10.1161/JAHA.125.046917}{Yang} & Short-term PM$_{2.5}$ and hospitalisation/ED use in atrial fibrillation & \textit{J Am Heart Assoc} & Case-crossover & PM$_{2.5}$ (wildfire-influenced) & USA\\
\addlinespace[0.9mm]
\multicolumn{6}{@{}l@{}}{\cellcolor{Mist}\textsf{\textbf{\textcolor{Sky}{Reproductive \& developmental}}}}\\[0.6mm]
\href{https://doi.org/10.1016/j.jes.2025.10.043}{Chen} & Prenatal air pollution, fetal growth trajectory phenotypes and child neurodevelopment & \textit{J Environ Sci (China)} & Prospective cohort & PM$_{2.5}$ / PM$_{10}$ & China\\
\href{https://doi.org/10.1016/j.ecoenv.2026.120566}{Huang} & Joint prenatal bisphenol + air-pollutant mixture exposure and preschool IQ & \textit{Ecotoxicol Environ Saf} & Prospective cohort (mixtures) & PM$_{2.5}$ / PM$_{10}$ & China\\
\href{https://doi.org/10.1007/s40572-026-00552-8}{Peterson} & Ambient air pollution, PAHs and pubertal development: critical review & \textit{Curr Environ Health Rep} & Critical review & PM$_{2.5}$ components & Global\\
\href{https://doi.org/10.1016/j.envres.2026.125313}{Xie} & Long-term PM$_{2.5}$ and timing of pubertal development in boys and girls & \textit{Environ Res} & Prospective cohort & PM$_{2.5}$ & China\\
\addlinespace[0.9mm]
\multicolumn{6}{@{}l@{}}{\cellcolor{Mist}\textsf{\textbf{\textcolor{Sage}{Respiratory \& allergic}}}}\\[0.6mm]
\href{https://doi.org/10.1007/s00484-026-03289-z}{Bayar Muluk} & Climate change and otolaryngology: airway disease impacts and emission reduction & \textit{Int J Biometeorol} & Narrative review & PM$_{2.5}$ & Global\\
\href{https://doi.org/10.4103/mgr.MEDGASRES-D-25-00164}{Liao} & Short-term air pollution, respiratory symptoms and lung function in asthma & \textit{Med Gas Res} & Panel study & PM$_{2.5}$ / PM$_{10}$ & China\\
\href{https://doi.org/10.3390/epidemiologia7040097}{Sanchez de Prada} & Ambient PM and asthma-coded admissions in children 0-3 y, Spain & \textit{Epidemiologia} & Ecological panel & PM$_{2.5}$ / PM$_{10}$ & Spain\\
\href{https://doi.org/10.1016/j.jhazmat.2026.143076}{Yang} & Allergenic vs. toxic risk of airborne seed hairs compared with ambient PM$_{2.5}$ & \textit{J Hazard Mater} & Field measurement & PM$_{2.5}$ + bioaerosol & China\\
\addlinespace[0.9mm]
\multicolumn{6}{@{}l@{}}{\cellcolor{Mist}\textsf{\textbf{\textcolor{Violet}{Mechanistic toxicology}}}}\\[0.6mm]
\href{https://doi.org/10.1016/j.jes.2025.12.074}{Ding} & PM$_{2.5}$ promotes vascular calcification via miR-27b-5p/HES1 signalling & \textit{J Environ Sci (China)} & In vitro + organoid + murine & PM$_{2.5}$ & China\\
\href{https://doi.org/10.1007/s00360-026-01698-5}{Maraschi} & Osmoregulatory disruption and metal bioaccumulation from settleable atmospheric PM & \textit{J Comp Physiol B} & Controlled exposure (ecotox) & Settleable APM & Brazil\\
\href{https://doi.org/10.1016/j.jes.2025.11.018}{Wu} & Macrophage NLRP3 inflammasome in real-ambient PM-induced hepatic glucose dysregulation & \textit{J Environ Sci (China)} & Real-ambient murine + co-culture & Real-ambient PM & China\\
\addlinespace[0.9mm]
\multicolumn{6}{@{}l@{}}{\cellcolor{Mist}\textsf{\textbf{\textcolor{Clay}{Occupational \& indoor}}}}\\[0.6mm]
\href{https://doi.org/10.3389/fmedt.2026.1796828}{Chen} & Nested suction instrument to cut intraoperative surgical-smoke PM$_{2.5}$ & \textit{Front Med Technol} & Bench / device experiment & PM$_{2.5}$ (surgical smoke) & China\\
\href{https://doi.org/10.1093/annweh/wxag055}{Cole} & Determinants of respirable crystalline silica exposure in Australian tunnelling & \textit{Ann Work Expo Health} & Observational exposure study & Respirable dust / RCS & Australia\\
\href{https://doi.org/10.1016/j.jhazmat.2026.143081}{Ding} & Stage-resolved hair Zn-isotope dynamics as a firefighting combustion-exposure tracer & \textit{J Hazard Mater} & Repeated-measures biomarker & Combustion-derived PM & China\\
\href{https://doi.org/10.3390/jox16040119}{Ramsden} & Xenobiotic hazards in aircraft cabin air & \textit{J Xenobiot} & Review + risk economics & Ultrafine particles & Global\\
\addlinespace[0.9mm]
\multicolumn{6}{@{}l@{}}{\cellcolor{Mist}\textsf{\textbf{\textcolor{Amber}{Sensing, forecasting \& instrumentation}}}}\\[0.6mm]
\href{https://doi.org/10.1016/j.jenvman.2026.130515}{Shi} & Adaptive spatiotemporal graph transformer for multi-site multi-step PM$_{2.5}$ forecasting & \textit{J Environ Manage} & Deep-learning model & PM$_{2.5}$ & China / India\\
\href{https://doi.org/10.3390/bios16070380}{Tsuchiya} & Seasonal and indoor-outdoor characteristics of fluorescent aerosol particles in offices & \textit{Biosensors} & Real-time sensor campaign & Bioaerosol / coarse & Japan\\
\addlinespace[0.9mm]
\multicolumn{6}{@{}l@{}}{\cellcolor{Mist}\textsf{\textbf{\textcolor{Amber}{Exposure assessment \& modelling}}}}\\[0.6mm]
\href{https://doi.org/10.1029/2025GH001621}{Solorzano-Araque} & Multiscale exposure modelling: LOTOS-EUROS + mobility surveys + particle toxicity & \textit{GeoHealth} & CTM + mobility integration & PM$_{2.5}$ / PM$_{10}$ & Colombia\\
\addlinespace[0.9mm]
\multicolumn{6}{@{}l@{}}{\cellcolor{Mist}\textsf{\textbf{\textcolor{Coral}{Burden, policy \& mitigation}}}}\\[0.6mm]
\href{https://doi.org/10.1002/hsr2.72886}{Karamzad} & Ischaemic heart disease burden in MENA 1990-2021; PM attributable fraction 28.2\% & \textit{Health Sci Rep} & GBD secondary analysis & Ambient PM & MENA\\
\href{https://doi.org/10.1371/journal.pone.0347989}{Rees} & HEPA filtration to reduce respiratory infection in care homes & \textit{PLoS One} & RCT process evaluation & Filtered PM & UK\\
\href{https://doi.org/10.1016/j.jes.2026.01.009}{Shan} & Policy-driven emission reductions and post-pandemic rebound, high-resolution inventories & \textit{J Environ Sci (China)} & Emission inventory & PM emissions & China\\
\addlinespace[0.9mm]
\multicolumn{6}{@{}l@{}}{\cellcolor{Mist}\textsf{\textbf{\textcolor{Slate}{Other clinical endpoints}}}}\\[0.6mm]
\href{https://doi.org/10.3390/dj14070440}{Lupita} & PM$_{2.5}$, preventive-care capacity and caries/edentulism burden across the EU & \textit{Dent J} & Ecological panel & PM$_{2.5}$ & EU-27\\
\href{https://doi.org/10.1016/j.jdin.2026.06.008}{Park} & Long-term PM exposure and risk of androgenetic alopecia & \textit{JAAD Int} & Nationwide cohort & PM$_{2.5}$ / PM$_{10}$ & South Korea\\
\href{https://doi.org/10.3171/2026.3.PEDS2615}{Sun} & Environmental exposure and area deprivation vs. paediatric hydrocephalus outcomes & \textit{J Neurosurg Pediatr} & Registry cohort + geospatial & PM$_{2.5}$ (satellite-derived) & USA\\
\addlinespace[0.9mm]
\bottomrule
\end{longtable}}


\vspace{4mm}
\begin{keybox}
{\sffamily\bfseries\footnotesize\color{Deep}Provenance and method}\\[1.2mm]
{\fontsize{7.2}{8.8}\selectfont
Entry window \textbf{3 September 2026}, single day, opened against a
\texttt{last\_entry\_date} of 2 September. Sources queried on 4 September 2026 at 11:3x
CDT; under this pipeline's 22:00-local rule, 3 September is the newest buildable date and
4 September is not opened. No day is owed: 1 and 2 September shipped on the previous run.
\textbf{PubMed} E-utilities on \texttt{[EDAT]} 2026/09/03: \textbf{29} hits on the health
axis, \textbf{6} on the instrumentation axis.
\textbf{Europe PMC} \texttt{CREATION\_DATE}: 23.
\textbf{Crossref} by-ISSN, windowed on \texttt{created}: \textbf{94} deposits.
\textbf{arXiv} answered normally this run (60 entries, no HTTP 429, the first clean
response since 31 August) but its newest submission is 2026-08-20, so the undated relevance
query again contributed nothing; \textbf{OpenAlex} returned HTTP 429 (the plan-gated date
filter has been broken since 28 July). \textbf{Consensus} returned ten low-cost-sensor
calibration papers, every one published 2018--2024 and therefore out of window --- the
fifth consecutive nil run for that leg. Full text was pulled through the browser for two
records only: the Plantower evaluation (Taylor \& Francis, JavaScript-gated) and the
Guangdong CHD paper, whose PubMed abstract renders its strongest component as
``[Formula: see text]'' and which was resolved to \emph{sulfate} at source rather than
guessed. ScienceDirect refused retrieval, which is why the two \textit{Atmos.\ Pollut.\
Res.} records ship on metadata only.
After collapsing cross-source duplicates by PMID, DOI and normalised title, \textbf{138
unique records} were considered; \textbf{12 were already shipped} in earlier issues,
\textbf{38 were screened in scope} and \textbf{88 were screened out}. Notable exclusions: a
tire-wear-particle earthworm toxicology paper and a Cu/SSZ-13 SCR catalyst study, both
airborne-adjacent but with no ambient PM metric or human endpoint; two \textit{Eur.\ J.\
Public Health} conference abstracts on indoor air quality in cystic fibrosis and the UKHSA
air-quality dashboard, both on topic but with no retrievable text; and a
\textit{Health Aff.\ Sch.} commentary on weather-resilient primary-care payment, which
carries no exposure metric.
\textbf{Trial watch:} zero registrations or updates for this date across five air-pollution
query terms. This is a true zero and not a query failure --- the same date filter with no
term returns 1{,}001 studies updated 3 September. Logged to \texttt{state/trials.json} for
the W36 weekly.
Two DOIs (\texttt{10.12182/20260760103}, \texttt{10.3967/bes2026.072}) are registered
outside Crossref and were initially failed as dead by the DOI gate; both resolve in the
Handle system and the gate was extended accordingly, so both links are live.
Subtopic, design, particle-metric and geography labels are single-label by dominant
contribution and are a judgement, not a controlled vocabulary. All DOI links resolve to the
publisher of record; abstracts remain the copyright of their respective publishers.\par}
\end{keybox}

\end{document}
